\Article{18}
When we apply the preceding investigation to the evaluation of the second term of the expression $\Omega$ in Article 3, the second region, denoted by $S$ from Article 6 onward, is identical with the first. What in Article 16 were $\sigma$, $\tau$, $\tau'$ will therefore here be $s$, $t$, $0$, if $t$ denotes the whole surface of the region $s$ filled with fluid. Thus whenever this region contains neither parts of perceptible extent but imperceptible thickness nor gaps (fissures) of this kind, the second part of the expression $\Omega$ becomes
\[
=\frac12\pi cc(s\psi0-t\theta0)
\]

There are accordingly two exceptions.

1) If the region $s$ contains a part of imperceptible thickness, its surface has two parts whose sizes do not differ perceptibly. Denoting either of them by $t'$ and the thickness of the region at any element $\mathrm{d}t'$ generally by $\rho$, we must add to the preceding expression the term
\[
\pi cc\int\theta'\rho.\mathrm{d}t'
\]

2) If the region $s$ contains a cavity of imperceptible thickness, a similar term is added, namely $\pi cc\int\theta'\rho.\mathrm{d}t''$, with $t''$ denoting either of the parts of the surface $t$ adjoining the fissure, and $\rho$ the thickness of the fissure at any point.

In evaluating the third term of the expression $\Omega$, we must retain the symbol $S$ to denote the region occupied by the vessel, but in place of the function symbol $f$ we must use $F$ \SourceLeaf{58}for the attractive force of the molecules of the vessel. Likewise, for the functions represented by the symbols $\varphi$, $\psi$, $\theta$, $\theta'$ we must use others represented by the symbols $\Phi$, $\Psi$, $\Theta$, $\Theta'$, which we suppose depend on $F$ in the same way as the former depend on $f$. What in the general investigation were $\sigma$, $\tau'$ will evidently be $0$ here. For $\tau$ we shall simply use the letter $T$ to denote only the part of the surface of the region $S$ that adjoins the fluid. In this way the third part of the expression $\Omega$ becomes, generally speaking,
\[
=\pi cCT\Theta0
\]
with two exceptions here as well, namely

3) If at a perceptible part $T'$ of the surface $T$ the fluid has an imperceptible thickness, to be denoted generally by $\rho$, the term to be added is
\[
-\pi cC\int\Theta'\rho.\mathrm{d}T'
\]

4) If, besides the part $T$ adjoining the fluid, the surface of the vessel has another part $T''$ at some distance from the fluid, albeit an imperceptible one, then, with $\rho$ denoting this distance generally for any point, the term to be added is
\[
+\pi cC\int\Theta'\rho.\mathrm{d}T''
\]
It would be superfluous to dwell on the first exception insofar as it is outside the third, or on the second or fourth. Even if equilibrium of the fluid could occur in certain cases of this kind, those cases would be very special, and such an equilibrium could be neither stable nor accessible to experiment. By contrast, the first exceptional case, insofar as it is included in the third, is indeed essential to the theory. We shall nonetheless set it aside for a while here so that we can investigate the conditions for equilibrium insofar as it can subsist without an imperceptible film of fluid adhering to the vessel.

Setting all these exceptions aside, then, the expression whose value must be a maximum in the state of equilibrium will be
\[
-gc\int z\mathrm{d}s+\frac12ccs\psi0-\frac12\pi cct\theta0+\pi cCT\Theta0
\]
and since the volume $s$ remains unchanged through every change the shape of the fluid can undergo, the following expression
\[
\int z\mathrm{d}s+\frac{\pi c\theta0}{2g}.t-\frac{\pi C\Theta0}{g}.T
\]
will have to be a \textit{minimum} in the state of equilibrium.

\SourceLeaf{59}We have already noted that $\frac{c\theta0}{g}$ represents a two-dimensional region, and the same holds for $\frac{C\Theta0}{g}$. If, then, we set
\[
\frac{\pi c\theta0}{2g}=\alpha\alpha,\qquad\frac{\pi C\Theta0}{2g}=\beta\beta
\]
$\alpha$, $\beta$ will be constant lengths depending on the relation of gravity to the strength of the forces that the parts of the fluid experience from one another and from the molecules of the vessel. If we further denote the free part of the surface of the fluid, that is, the part that does not adjoin the vessel, by $U$, so that $t=T+U$, the following expression, henceforth denoted by $W$, must be a minimum in the state of equilibrium
\[
\int z\mathrm{d}s+(\alpha\alpha-2\beta\beta)T+\alpha\alpha U
\]

\Article{19}
Before we work out in full generality what follows from this important theorem, it will be worthwhile to show how readily it yields the principal phenomenon of capillary tubes.

Let us consider a fluid in equilibrium in a vessel with two limbs, with one part of the free surface of the fluid in the first limb and another part in the second. We suppose that the walls of the vessel are vertical where they border these parts. Let $a$ be the area of the internal horizontal cross-section of the first limb (or, more precisely, of the horizontal projection of the free surface of the fluid in the first limb), $b$ its perimeter, and finally $ah$ the volume of fluid in this limb with the vertical wall extended downward to the plane from which the distances $z$ are measured. Equivalently, let $h$ be the mean height of the fluid above this plane. Let $a'$, $b'$, $h'$ denote the corresponding quantities for the second limb. If we imagine that the state of the fluid undergoes an infinitely small change in which both parts of the free surface retain their shape, the variation of the first part of the expression $W$, namely the integral $\int z\mathrm{d}s$, will evidently be
\[
=ah\mathrm{d}h+a'h'\mathrm{d}h
\]
and the variation of $T$ will be
\[
=b\mathrm{d}h+b'\mathrm{d}h'
\]
and finally by hypothesis $\mathrm{d}U=0$. From this we obtain
\[
\mathrm{d}W=ah\mathrm{d}h+a'h'\mathrm{d}h'-(2\beta\beta-\alpha\alpha)(b\mathrm{d}h+b'\mathrm{d}h')
\]
\SourceLeaf{60}Further, since the total volume of the fluid remains unchanged, we have
\[
a\mathrm{d}h+a'\mathrm{d}h'=0
\]
and hence
\[
\mathrm{d}W=\mathrm{d}h[a(h-h')-(2\beta\beta-\alpha\alpha)(b-\frac{ab'}{a'})]
\]
The condition that $W$ be a minimum in the state of equilibrium thus leads to the equation embodying the principal phenomenon of capillary tubes
\[
h-h'=(2\beta\beta-\alpha\alpha)(\frac{b}{a}-\frac{b'}{a'})
\]
and it is immediately evident that this equation does indeed correspond to a \textit{minimum} of $W$, since the value of $\frac{\mathrm{d}\mathrm{d}W}{\mathrm{d}h^2}$ is $=a+\frac{aa}{a'}$, which is positive by its nature.

We call the second limb wider than the first whenever $\frac{a'}{b'}$ is greater than $\frac{a}{b}$. The fluid will therefore be depressed or raised more in the narrower limb than in the wider one according as the square $\beta\beta$ is less or greater than $\frac12\alpha\alpha$. If it should happen that $\beta\beta=\frac12\alpha\alpha$, the height would be the same in both limbs. If the second limb is so wide that $\frac{b'}{a'}$ can be neglected compared with $\frac{b}{a}$, we have approximately
\[
h-h'=(2\beta\beta-\alpha\alpha)\frac{b}{a}
\]
Thus in cylindrical capillary tubes the depression or rise of the fluid is inversely proportional to the diameter of the tube. All of this agrees both with experiment and with what the illustrious \textsc{Laplace} sought to establish by theory.

If the vessel has several vertical limbs communicating with one another, let $a''$, $b''$, $h''$ for the third limb, $a'''$, $b'''$, $h'''$ for the fourth, and so on, denote what $a$, $b$, $h$ denote for the first. Then we also have
\[
\begin{gathered}
h-h''=(2\beta\beta-\alpha\alpha)(\frac{b}{a}-\frac{b''}{a'})\\
h-h'''=(2\beta\beta-\alpha\alpha)(\frac{b}{a}-\frac{b'''}{a'''})
\end{gathered}
\]
These equations can be written more neatly as
\[
\begin{aligned}
h-(2\beta\beta-\alpha\alpha)\frac{b}{a}&=h'-(2\beta\beta-\alpha\alpha)\frac{b'}{a'}\\
&=h''-(2\beta\beta-\alpha\alpha)\frac{b''}{a''}=h'''-(2\beta\beta-\alpha\alpha)\frac{b'''}{a'''}\text{etc.}
\end{aligned}
\]
Since the horizontal plane from which heights are measured is arbitrary, if we choose it so that
\[
h=(2\beta\beta-\alpha\alpha)\frac{b}{a}
\]
\SourceLeaf{61}it is evident that in the other limbs too we shall have
\[
\begin{gathered}
h'=(2\beta\beta-\alpha\alpha)\frac{b'}{a'},\qquad h''=(2\beta\beta-\alpha\alpha)\frac{b''}{a''},\\
h'''=(2\beta\beta-\alpha\alpha)\frac{b'''}{a'''}\quad\text{etc.}
\end{gathered}
\]
This plane, whose definition we shall establish more generally below, may be called the normal horizontal plane (plan de niveau). Supposing the vertical walls of the individual limbs to be extended to this plane if necessary, $ah$, $a'h'$, $a''h''$ etc. will express, for $2\beta\beta>\alpha\alpha$, the quantities of fluid raised above this plane in the individual limbs, or for $2\beta\beta<\alpha\alpha$, the quantities of fluid missing below this plane in the individual limbs. These quantities are thus equal to the products of the constant area $2\beta\beta-\alpha\alpha$ and the perimeters $b$, $b'$, $b''$, $b'''$ etc.

